\begin{abstract}
A number on a dashboard is a claim about a chain of SQL, orchestration code and business intent
that nobody can read end to end. Lineage tools draw that chain but assert it; they do not check it,
and they rarely say how a value is calculated or what it means. Konstruct is an open-source,
local-first system that reads a data repository and its warehouse (dbt, SQL, Airflow, Dagster,
PySpark, Looker, Power BI, Fivetran and Airbyte) and builds one graph of every relation, column,
task, test, dashboard and business term, with column-level lineage parsed statically from the SQL
that connects them. It composes each column's expression through every model upstream into a
formula over source columns, derives an ontology of entities, measures and relationships from the
code and its tests, reads and writes Apache Ossie semantic models, and gives coding agents the
graph through an MCP server and a gate that a change must pass. Its central idea is that lineage
should be measured, not asserted: \cmd{konstruct verify} perturbs each input column of each
model, reruns the model's SQL and compares what changed with what the graph claims, so that a
missing edge shows up as a dependency the graph does not explain, reproducible from the recorded
seed. On dbt Labs' jaffle-shop all \jsConfirmed{} of \jsClaimed{} claimed column edges are
confirmed with \jsUnexplained{} unexplained influences; on the bundled example, of
\exDeletedEdges{} verified edges whose input is not also a join, filter or grouping column, each
deleted in turn, verification reports \exDeletedEdgesPct{}\% as unexplained; \ossieNamedPortable{}
of the \ossieNamedElements{} elements of the Apache Ossie portability fixtures in a form the
specification allows survive a round trip unchanged; and a cold build of \perfLargestModels{}
generated models with \perfLargestEdges{} column edges takes \perfLargestCold{}~s on
\machineCores{} CPUs. Every measurement of Konstruct in this paper is regenerated from the code by
continuous integration.
\end{abstract}
